\section{IPV Restriction and Pharmaceutical Procurement}\label{app:robustness_ipv}

A reader might argue that pharmaceutical procurement involves item-level specifications---ANVISA compliance, brand-name regulatory history, storage chain, production line provenance---that are not well captured by an independent-private-values model with only type-level asymmetry. The concern is that what looks like a bidder-specific private value is in fact a common quality signal that correlates bidder valuations within an auction, in which case Assumption~\ref{a:ipv} is misspecified and the estimated $F_c^k$ absorbs variation that should be attributed to a common factor.

Three features of the BEC setting address this concern. First, the sample is restricted to Preg\~ao, the electronic reverse-auction format reserved for goods with standardized specifications at the item level; BEC's CADMAT catalog is designed to strip cross-item quality variation within class. This is a different data-generating environment from the construction contracts studied in \citet{bajari2014} or the Bayesian auctions of \citet{hendricks2003}, where contract-specific complexity drives correlation. Second, the CMED price ceiling applied to pharmaceutical class 6531 is itself a standardization device: any CMED-regulated drug clears a uniform quality bar, and the residual variation across suppliers is on production cost and distribution scale rather than on quality primitives. Third, the Krasnokutskaya auction-level heterogeneity correction (Assumption~\ref{a:uh}) allows for an auction-common component $a_t$ that absorbs exactly the kind of unobserved scale or specification variation the objection raises. The Turnbull NPMLE robustness relaxes Assumption~\ref{a:uh} entirely and delivers the same decomposition within 16 percent of the total effect across classes---with the largest gap (16 percent) in pharmaceutical Turnbull and a 5 percent gap in non-pharmaceutical Turnbull (Table~\ref{tab:v3_sensitivity_fc})---so the main results do not hinge on the Gaussianity of $a_t$ or on the correctness of the BLP shrinkage.

The pharmaceutical welfare-dominance result of Section~\ref{sec:discussion} is therefore subject to two qualifications. First, the strict-invariance check of Section~\ref{sec:robustness} reverses the pharmaceutical ranking: the welfare-dominance result in pharmaceuticals is conditional on the equilibrium-selection treatment of $F_c^{\text{SME,Post}}$. Second, the IPV diagnostic supports Assumption~\ref{a:ipv} on the available evidence but does not constitute a definitive test. The non-pharmaceutical welfare-dominance result is preserved across both qualifications.

\section{Coincident Regulatory Events at the March 2018 Cutoff}\label{app:cutoff}

The March 2018 cutoff might in principle coincide with other Brazilian regulatory events that could confound identification---CMED adjustments, platform upgrades, tax changes, federal procurement reform. Each is worth checking explicitly.

CMED list-price adjustments occur annually at the beginning of each year and apply to all CMED-regulated drugs, not to Group 65 specifically; a CMED adjustment would appear in the pre-period Group-65 data as well as in the never-treated product groups of the DiD benchmark (Appendix~\ref{app:did}), and the DiD would difference it out. BEC platform architecture was stable throughout the sample window; the catalog and auction format did not change. Federal Law 14{,}133/2021 did not exist during the window---the relevant federal framework was the 2006 statute, unchanged throughout. Any Group-65-specific shock that could substitute for the SME rule would need to coincide with March 2018, hit only Group 65, and leave no trace in publicly available administrative records. The DiD pre-trend and placebo diagnostics on the same sample are consistent with the timing being uncorrelated with such shocks; the present design treats the cutoff as plausibly exogenous on those diagnostics rather than as definitively so.
